% chapters/memory.tex — Memory chapter (narrative text only)
% BiDi host chapter: contains the substantive Hebrew/English section.
\chapter{Memory and State Persistence}
\label{chap:memory}

Long-horizon agentic tasks require agents to retain information across many
steps and interaction turns.
Modern \LLM{}-based agents address this by maintaining explicit memory
stores alongside the model's implicit parametric knowledge~\cite{huang2025licomemory,chen2025telemem}.
\index{memory}
\index{state persistence}

\section{Memory Taxonomies}
\label{sec:memory-taxonomy}

Agent memory is commonly classified along two axes:

\begin{itemize}
  \item \textbf{Episodic memory} stores raw experience traces (observations,
        actions, rewards) and supports retrieval by similarity or recency.
  \item \textbf{Semantic memory} stores distilled facts and concepts
        extracted from experience, enabling efficient generalisation.
  \item \textbf{Working memory} holds the current context window of the
        \LLM{} and is the primary bottleneck for long-horizon tasks.
\end{itemize}

\section{Lightweight Cognitive Architectures}
\label{sec:memory-lightcog}

LiCoMemory~\cite{huang2025licomemory} introduces a lightweight cognitive
memory system that reduces storage overhead while preserving recall accuracy
on long-reasoning benchmarks.
TeleMem~\cite{chen2025telemem} extends this to multimodal long-term memory,
maintaining both textual and visual traces across extended interaction histories.

\section{Memory Mechanisms in Hebrew / מנגנוני זיכרון בעברית}
\label{sec:memory-bidi}

This section demonstrates Hebrew/English bidirectional typesetting.
The body text is Hebrew (right-to-left); embedded English technical terms
switch direction inline via \texttt{polyglossia}'s \texttt{\textbackslash textenglish} command.
\index{bidirectional text}
\index{זיכרון@\texthebrew{זיכרון} (memory)}

\medskip
\begin{hebrew}
מנגנוני זיכרון הם מרכיב מרכזי בכל מערכת \textenglish{LLM} אגנטית.
הסוכן (\textenglish{agent}) \texthebrew{חייב לשמור מידע}
לאורך פרקי זמן ארוכים כדי לבצע משימות מורכבות.

גישות כמו \textenglish{LiCoMemory}~\textenglish{\cite{huang2025licomemory}}
ו-\textenglish{TeleMem}~\textenglish{\cite{chen2025telemem}}
מציגות ארכיטקטורות זיכרון קלות משקל
שמגדילות את יכולת הסוכן לשמור מידע ולאחזר אותו בצורה יעילה.
שיטות אלו מאזנות בין עלות האחסון (%
\textenglish{storage cost}) לבין דיוק האחזור (\textenglish{recall accuracy}).
\end{hebrew}
\medskip

The paragraph above is typeset in Hebrew with correct right-to-left alignment
and bidirectional text flow.  English citations and terms appear left-to-right
as inline switches, satisfying the PRD~§8.6 BiDi requirement.
